\documentclass{exam-zh}

\begin{document}
\chapter*{Issue \#30 验证文档}

\section*{问题与验证范围}

旧实现会先捕获整个 \texttt{choices} 环境，再把每个选项直接放入盒子测宽。
因此原始复现中的 \verb|\verb| 会在 \verb|\end{choices}| 处触发编译错误。
当前实现只对含 \verb|\verb| 的选项重新扫描，然后继续使用原来的自动测宽、
多列排版和正确答案标记逻辑。

\section*{修复结果}

\begin{question}
  以下程序的运行结果为 \paren。

  \begin{verbatim}
#include <iostream>
int main() {
  std::cout << "Hello, world!" << std::endl;
}
  \end{verbatim}

  \begin{choices}[columns=2,show-answer=true]
    \item \verb+Alpha+
    \item* \verb-Beta-
    \item \verb,Gamma,
    \item \verb|Delta|
  \end{choices}
\end{question}

上面的四个选项能够正常编译并保持两列布局；正确选项 \verb|Beta| 仍按原有规则着色。

\section*{当前边界}

\verb|question| 和 \verb|problem| 可直接放置 \verb|verbatim|；
\verb|choices| 现已支持本例中的内联 \verb|\verb|。
直接嵌入 \verb|lstlisting| 仍受整段正文捕获和换行 catcode 限制，
多行代码选项应继续使用 \verb|\lstinputlisting| 从外部文件读取。

\end{document}
